% !TEX root = ../main.tex

\chapter{模 \texorpdfstring{$\ell^n$}{l n} 与模 \texorpdfstring{$p$}{p} 的 \texorpdfstring{$L$}{L}-函数}\label{III}

在本讲中, 我对与 $p$ 互素的挠系数环 $A$, 或特征为 $p$ 的域 $A$, 证明一个与迹公式报告 (本书, 以下简称“报告”) 的主定理类似的公式. 证明将实质性地使用 \cite[XVII 5.5]{sga4}.

\section{对称张量}\label{III:1}

\subsection{}\label{III:1-1}

设 $A$ 是交换环. 函子 $\Gamma^n:(\text{$A$-模})\to (\text{$A$-模})$ 的性质列表参见 \cite[XVII 5.5.1 and 2]{sga4}. 我们只回顾, 当 $M$ 平坦时 (实际上我们只关心 $M$ 为有限型投射模的情形), 有
\[
  \Gamma^n M = \left(M^{\otimes n}\right)^{S_n} \qquad \text{(次数为 $n$ 的对称张量).}
\]

\subsection{}\label{III:1-2}

设 $X$ 是 $S$-概形, $X^n/S$ 是 $X$ 的第 $n$ 个纤维幂, $\pi$ 是从 $X^n/S$ 到 $(X^n/S)/S_n=\operatorname{Sym}_S^n(X)$ 的投影 (假设这个商存在). 若 $\sG$ 是 $X$ 上的 $A$-模层, 则 $\sG$ 的 $n$ 重外张量积 $\sG^{\boxtimes n}$ 是 $X^n/S$ 上的 $S_n$-等变层. 我们将用到 \cite[XVII 5.5.7 to 9]{sga4} 的\emph{外}函子 $\Gamma^n$. 只回顾, 当 $\sG$ 为平坦 $A$-模层时 (实际上我们只关心可构造平坦的情形), 有
\[
  \Gamma_\text{ext}^n(\sG) = \left(\pi_*\sG^{\boxtimes n}\right)^{S_n} \text{.}
\]

\begin{lemma_}\label{III:1-3}
设 $S=\spec(k)$, 其中 $k$ 代数闭, $X$ 是有限个 $S$ 的副本的不交并, $G$ 是 $X$ 上的平坦 $A$-模层, 则有
\[
  \Gamma\left(\operatorname{Sym}_S^n(X),\Gamma_\text{ext}(\sG)\right) = \Gamma^n \Gamma(X,\sG)
\]
\end{lemma_}

事实上,
\begin{align*}
  \Gamma\left(\operatorname{Sym}_S^n(X),\Gamma_\text{ext}^n(\sG)\right)
    &= \Gamma\left(\operatorname{Sym}_S^n(X),\pi_*\sG^{\boxtimes n}\right)^{S_n} \\
    &= \Gamma\left(X^n/S,\sG^{\boxtimes n}\right)^{S_n}
    = \left(\Gamma(X,\sG)^{\otimes n}\right)^{S_n}
    = \Gamma^n \Gamma(X,\sG) \text{.}
\end{align*}

\subsection{}\label{III:1-4}

我们将用到非加性函子 $\Gamma^n$ 与 $\Gamma_\text{ext}^n$ 的导出函子.

这类导出函子的理论由 Dold 与 Puppe 建立 \cite{do61}. 比下面的回顾更完整的说明参见 \cite[XVII 5.5.3]{sga4}.

设 $\cA$ 是阿贝尔范畴 (甚至只须是其中每个幂等自同态都是到某个直和因子上的投影的加性范畴). 以 $N$ 表示正规化函子: $N:($由 $\cA$ 的对象组成的余单纯复形$)\to($由 $\cA$ 的对象组成的微分复形, 当 $i<0$ 时为 $K^i)$ . 它是范畴等价. 它及其逆都将同伦态射送到同伦态射. 若 $\Gamma$ 是函子 $\cA\to\cB$, $K$ 是 $\cA$ 的对象组成的非负次数复形 ($K^i=0$ 对 $i<0$), 则置 $T K=N T N^{-1} K$.

\subsection{例}\label{III:1-5}

若 $K$ 在次数 $0$ 集中为 $M$, 则 $T K$ 在次数 $0$ 集中为 $T M$.

\subsection{}\label{III:1-6}

若 $K$ 是有限型投射 $A$-模组成的非负次数有界复形, 则 $A$-模复形 $\Gamma^n K$ 具有同样的性质.

用 $\eD_\text{parf}^{\geqslant 0}(A)$ 表示 $\eD_\text{parf}(A)$ (报告, \ref{II:4-3}) 中由 Tor 维数 $\leqslant 0$ 的复形构成的子范畴. 可以看出 $\Gamma^n$ 诱导函子
\[
  \eL\Gamma^n : \eD_\text{parf}^{\geqslant 0}(A) \to \eD_\text{parf}^{\geqslant 0}(A) \text{.}
\]

$\eL\Gamma_\text{ext}^n$ 的定义类似但更复杂, 参见 \cite[XVII 5.5.14]{sga4}.

\subsection{}\label{III:1-7}

设 $K$ 是有限型投射 $A$-模的有界复形, $u$ 是 $K$ 的自同态. 用 $\det(1-u t,K)$ 表示常数项为 $1$ 的形式级数
\[
  \det(1-u t,K) = \prod_i \det(1-u t,K^i)^{(-1)^i} \text{.}
\]

下列性质的验证留给读者.

\subsubsection{}\label{III:1-7-1}

设 $F$ 是在 $u$ 下稳定的有限滤过, 且各 $\gr_F^p(K^q)$ 都是有限型投射模. 则
\[
  \det(1-u t,K) = \prod_p \det(1-u t,\gr_F^p(K)) \text{.}
\]

\subsubsection{}\label{III:1-7-2}

将 $K$ 向左平移 $q$ 位, 则有
\[
  \det(1-u t,K[q]) = \det(1-u t,K)^{(-1)^q} \text{.}
\]

\begin{proposition_}\label{III:1-8}
设 $u$ 是有限型投射 $A$-模组成的非负次数有界复形的自同态. 有 \
\begin{equation*}\tag{1.8.1}\label{III:eq:1-8-1}
  \det(1-u t,K)^{-1} = \sum_n \tr(u,\Gamma^n K) t^n \text{.}
\end{equation*}
\end{proposition_}

将 \eqref{III:eq:1-8-1} 的右边记为 $D(u,K)$.

\begin{lemma_}\label{III:1-9}
设 $u$ 是有限型投射 $A$-模组成的次数 $\geqslant 0$ 的有界复形短正合列
\[
  0 \to K' \to K \to K'' \to 0 \text{.}
\]
的自同态. 则有 $D(u,K)=D(u,K')\cdot D(u,K'')$.
\end{lemma_}

若 $0\to M'\to M\to M''\to 0$ 是 $A$-模的分裂短正合列, 则 $\Gamma^n M$ 有自然滤过, 其逐次商为各 $\Gamma^i M'\otimes \Gamma^j M''$ ($i+j=n$). 因此, $\Gamma^n K$ 有自然滤过, 其逐次商为复形 $N(N^{-1} \Gamma^i K'\otimes N^{-1}\Gamma^j K'')$ ($i+j=n$), 它们典范地同伦于双复形 $\Gamma^i K'\otimes \Gamma^i K''$ 所对应的单复形 $s(\Gamma^i K'\otimes \Gamma^i K'')$, 且
\[
  \tr(u,\Gamma^n K) = \sum_{i+j=n} \tr(u,\Gamma^i K')\cdot \tr(u,\Gamma^j K'') \text{,}
\]
从而引理成立. 由 \ref{III:1-9} 归纳得到下述结果.

\begin{lemma_}\label{III:1-10}
设 $F$ 是 $K$ 的在 $u$ 下稳定的有限滤过, 各 $\gr_F^p(K^q)$ 为有限型投射模. 则
\[
  D(u,K) = \prod_p D\left(u,\gr_F^p(K)\right) \text{.}
\]
\end{lemma_}

\begin{lemma_}\label{III:1-11}
设 $u$ 是有限型投射 $A$-模 $M$ 的自同态, $M[-q]$ 是在次数 $q$ 集中为 $M$ 的复形. 有
\[
  D\left(u,M[-q]\right) = \left(\sum \tr(u,\Gamma^n M) t^n\right)^{(-1)^q} \qquad \text{($q\geqslant 0$.)}
\]
\end{lemma_}

由 \ref{III:1-5}, 当 $q=0$ 时成立. 用归纳法, 只须证明
\[
  D\left(u,M[-q]\right) \cdot D\left(u,M[-(q+1)]\right) = 1 \text{.}
\]
将 \ref{III:1-9} 应用于在次数 $q$ 与 $q+1$ 集中为 $M$ 的复形 $K$ (其 $d=\operatorname{id}$), 以及它的平凡滤过:
\[
  (\cdots 0 \to M \cdots) \to (\cdots M \to M) \to (\cdots M \to 0 \to 0 \cdots) \text{.}
\]
即得此公式. 事实上, $K$ 同伦于 $0$. 因而在导出范畴中有 $\Gamma^n(K)=\Gamma^n(0)$: 当 $n>0$ 时为 $0$, 当 $n=0$ 时为 $A$, 且 $D(u,K)=1$.

\subsection{}\label{III:1-12}

现在证明 \ref{III:1-8}. 将 \ref{III:1-10} 应用于 $K$ 的由子复形
\[
  0 \to \cdots \to 0 \to K^q \to K^{q+1} \to \cdots \text{,}
\]
组成的“平凡”滤过, 得到
\[
  D(u,K) = \prod_q D\left(u,K^q[-q]\right) =_{\text{(\ref{III:1-10})}} \prod_q \left(\sum \tr\left(u,\Gamma^n(K^q)\right) t^n \right)^{(-1)^q} \text{.}
\]
剩下只须证明, 对有限型投射模 $M$ 的自同态 $u$, 有
\begin{equation*}\tag{1.12.1}\label{III:eq:1-12-1}
  \det\left(1-u t,M\right)^{-1} = \sum \tr(u,\Gamma^n M) t^n \text{.}
\end{equation*}
若 $u=0$, 则两边均为 $1$. 在 $M$ 上添加一个其上 $u=0$ 的模 $N$, 再应用 \ref{III:1-9}, 即可归结为假设 $M$ 自由. 选取一个基. 此时须证明关于 $u$ 的坐标的一些代数恒等式. 根据代数恒等式的扩张原理, 可归结为假设 $A$ 是代数闭域. 模 $M$ 于是有在 $u$ 下稳定且逐次商的秩为 $1$ 的滤过, 再由 \ref{III:1-10} 归结到 $M$ 的秩为 $1$ 的情形. 此时公式 \eqref{III:eq:1-12-1} 化为恒等式
\[
  \frac{1}{1-a t} = \sum a^n t^n \text{.}
\]

\begin{corollary_}\label{III:1-13}
设 $u$ 和 $u'$ 分别为有限型投射 $A$-模的有界复形 $K$ 与 $K'$ 的自同态. 若在导出范畴中, 带有 $u$ 的 $K$ 同构于带有 $u'$ 的 $K'$, 则
\[
  \det(1-u t,K) = \det(1-u' t,K') \text{.}
\]
\end{corollary_}

应用 \ref{III:1-7-2}, 可归结为假设 $K$ 与 $K'$ 都在非负次数. 此时注意 \ref{III:1-8} 的右边具有所要的不变性即可.

\subsection{}\label{III:1-14}

此推论使我们可以对 $K\in \ob\eD_\text{parf}(A)$ 的自同态 $u$ 定义 $\det(1-u t,K)$. 这是一种针对当前问题的构造; 事实上, 对 $L\in\ob\eD_\text{parf}(\Lambda)$ 的自同构 $v$, 可以定义 $\det(v)\in\Lambda^\times$, 而取 $v=1-u t$, $\Lambda=A\llbracket t\rrbracket$, $L=K\otimes_A\Lambda$, 即得到 $\det(1-u t,K)$.

\section{定理}\label{III:2}

\subsection{}\label{III:2-1}

沿用 (报告, \S\ref{II:1}) 的记号. 对交换 Noether 挠环 $A$, $\dF_q$ 上的分离有限型概形 $X_0$, 以及 $\sK_0\in\ob\eD_\text{ctf}(X_0;A)$, 置
\[
  L(X_0,\sK_0) = \prod_{x\in |X_0|} \det\left(1-F_x^* t^{\deg(x)}, \sK\right)^{-1}
\]
(\ref{III:1-7}, \ref{III:1-13}, 以及报告, \ref{II:1-5}, \ref{II:1-6}).

\begin{theorem_}\label{III:2-2}
在以下两种情形的每一种中
\begin{enumerate}[\indent a)]
  \item $A$ 是与 $p$ 互素的挠环;
  \item $A$ 既约且特征为 $p$
\end{enumerate}
都有
\begin{equation*}\tag{2.2.1}\label{III:eq:2-2-1}
  L(X_0,\sK_0) = \det\left(1-F^* t^f, \eR\Gamma_c(X,\sK)\right)^{-1} \text{.}
\end{equation*}
\end{theorem_}

\subsection{注}\label{III:2-3}

在情形 a) 中, 利用报告 \S\ref{II:3} 的方法, 可从迹公式 (报告 \ref{II:4-10}) 推出 \eqref{III:eq:2-2-1} 两边的对数导数相等. 由于环 $A$ 是挠环, 这不足以证明 \ref{III:2-2}.

\subsection{注}\label{III:2-4}

若 $A$ 是域, 则有
\[
  L(X_0,\sK_0) = \prod_i L\left(X_0,\underline\h^i(\sK_0)\right)^{(-1)^i}
\]
而谱序列 $E_2^{p q} = \h_c^p(X,\underline\h^q(\sK))\Rightarrow \h_c^{p+q}(X,\sK)$ 说明
\begin{align*}
  \det\left(1-F^* t, \eR\Gamma_c(X,\sK)\right)
    &= \prod_n \det\left(1-F^* t^f,\h_c^n(X,\sK)\right)^{(-1)^n} \\
    &= \prod_{p,q} \det\left(1-F^* t,\h_c^p(X,\underline\h^q(\sK))\right)^{(-1)^{p+1}} \text{.}
\end{align*}
这将 \ref{III:2-2} (当 $A$ 为域时) 归结为以下特例.

\subsubsection{}\label{III:2-4-1}

若 $\sG_0$ 是可构造 $A$-向量空间层, 则有
\[
  L(X_0,\sG_0) = \prod_n \det\left(1-F^* t^f,\h_c^n(X,\sG)\right)^{(-1)^{n+1}} \text{.}
\]

\subsection{注}\label{III:2-5}

容易将情形 b) 归结到 $A$ 是特征为 $p$ 的域 (甚至有限域) 的情形; 参见 \ref{III:4-2}.

\subsection{注}\label{III:2-6}

b) 中 $A=\dZ/p$, 且 $X$ 化为在次数 $0$ 的常值层 $\dZ/p$ 的特例,
\[
  Z(X_0,t) = \prod_{x\in |X_0|} \left(1-t^{\deg(x)}\right)^{-1}
      \equiv \prod_n \det\left(1-F^* t^f,\h_c^n(X,\dZ/p)\right)^{(-1)^{n+1}} \pmod p
\]
已由 N. Katz \cite[XXII 3.1]{sga7} 证明. 他的方法与这里的方法完全不同, 从 Dwork 的 $Z(X_0,t)$ 的 $p$-进理论出发.

\subsection{注}\label{III:2-7}

与我曾经轻率地向朋友们断言的相反, 在 b) 中不能去掉“$A$ 既约”这一假设. 反例参见 \ref{III:4-5}.

\subsection{\ref{III:2-2} 的证明概要}\label{III:2-8}

通过次数平移 \ref{III:1-7-2}, 可归结为假设 $\sK_0$ 是可构造平坦 $A$-模层组成的非负次数有界复形. 此时有 $\eR\Gamma_c(X,\sK)\in \ob\eD_\text{parf}^{\geqslant 0}(A)$ (报告, \ref{II:4-9} 的证明 b). 我们将用 \ref{III:1-8} 把 \eqref{III:eq:2-2-1} 两边展开成关于 $T=t^f$ 的幂级数. $T^n$ 的系数相等, 将由 \cite[XVII 5.5.21]{sga4} 以及 $\operatorname{Sym}^n(X)$ 上的迹公式得到: 在情形 a) 中用 (报告, \ref{II:4-10}) 的公式, 在情形 b) 中用下一章将证明的公式. 首先须利用 \eqref{III:eq:2-2-1} 两边关于 $X_0$ 的乘法性, 归结到 $X_0$ 的对称幂存在的情形 (例如 $X$ 仿射的情形): 对 $X_0$ 的开子概形 $U_0$, 其闭补集为 $Y_0$, 有 $L(X_0,\sK_0)=L(U_0,\sK_0)\cdot L(Y_0,\sK_0)$, 而右边的类似公式由 \ref{III:1-7-1}, 按 (报告 \S\ref{II:6}A) 的方法得到.

\subsection{$X_0$ 有限的情形}\label{III:2-9}

利用 \eqref{III:eq:2-2-1} 两边关于 $X_0$ 的乘法性, 这一情形归结到 $X_0=\spec(\dF_{q^k})$. 回到定义可见, 此时 \eqref{III:eq:2-2-1} 等价于报告 \ref{II:3-4}.

\subsection{左边}\label{III:2-10}

由 \ref{III:1-8}, 有
\[
  L(X_0,\sK_0) = \prod_{x\in |X_0|} \sum_n \tr\left(F_x^*,\Gamma^n\sK\right) t^{m\deg(x)} \text{.}
\]
因此, $T^n$ 的系数是对 $|X_0|$ 中的元素的形式线性组合 $\sum_{x\in |X_0|} n_x\cdot x$ ($n_x\geqslant 0$, 且几乎所有 $n_x$ 为零) 中那些满足 $\sum n_x [k(x):\dF_q] = n$ 的组合, 将各乘积
\begin{equation*}\tag{2.10.1}\label{III:eq:2-10-1}
  \prod_x \tr\left(F_x^*,\Gamma^{n_x} \sK\right) \text{.}
\end{equation*}
相加所得.

对 $x\in |X_0|$, 用 $(x)$ 表示由 $X$ 中位于 $x$ 之上的 $[k(x):\dF_q]$ 个点组成的 $0$-闭链 (即 $F$ 的一个轨道). 映射 $\sum n_x\mapsto \sum n_x\cdot (x)$ 将上面满足 $\sum n_x [k(x):\dF_q]=n$ 的形式线性组合, 与 $X$ 上在 $F$ 下不变的次数为 $n$ 的有效 $0$-闭链等同, 也即与 $F$ 在 $\operatorname{Sym}^n(X)$ 中的不动点等同. 因此, $L(X_0,\sK_0)$ 中 $T^n$ 的系数表现为以 $\operatorname{Sym}^n(X)^F$ 为指标的各乘积 \eqref{III:eq:2-10-1} 的和. 更确切地说,

\begin{lemma_}\label{III:2-11}
 $L(X_0,\sK_0) = \sum_n \sum_{x\in\operatorname{Sym}^n(X)^F} \tr\left(F_x^*,\eL\Gamma_\textnormal{ext}^n \sK\right) T^n$.
\end{lemma_}

随着 $X_0$ 的有限子概形 $Y_0$ 越来越大, 有
\[
  L(X_0,\sK_0) = \lim L(Y_0,\sK_0) \text{,}
\]
即对每个 $m$, 都存在有限子概形 $Z_0\subset X_0$, 使得若 $Y_0\supset Z_0$, 则 $L(X_0,\sK_0)$ 与 $L(Y_0,\sK_0|Y_0)$ 模 $T^m$ 同余.
只须 $Z_0$ 包含所有次数 $\leqslant m$ 的闭点即可. 由于右边也有同样的性质, 只须在 $X_0$ 有限时证明 \ref{III:2-11}. 此时有
\begin{align*}
  L(X_0,\sK_0)
    &=_\text{\ref{III:2-9}} \det\left(1-F^* T,\Gamma(X,\sK)\right)^{-1} \\
    &=_\text{\ref{III:1-8}} \sum_n \tr\left(F^*,\Gamma^n \Gamma(X,\sK)\right) T^n \\
    &=_\text{\ref{III:1-3}} \sum_n \tr\left(F^*,\Gamma(\operatorname{Sym}^n(X),\Gamma_\text{ext}^n(\sK)\right) T^n \\
    &= \sum_n \sum_{x\in \operatorname{Sym}^n(X)^F} \tr\left(F_x^*,\Gamma_\text{ext}^n(\sK)\right) T^n
\end{align*}
(利用有限集这一平凡情形的迹公式).

\subsection{右边}\label{III:2-12}

由 \ref{III:1-8} 及 \cite[XVII 5.5.12]{sga4}, 有
\begin{align*}
  \det\left(1-F^* T,\eR\Gamma_c(X,\sK)\right)^{-1}
    &= \sum_n \tr\left(F^*,\eL\Gamma^n\eR\Gamma_c(X,\sK)\right) T^n \\
    &= \sum_n \tr\left(F^*,\eR\Gamma_c(\operatorname{Sym}^n(X),\Gamma_\text{ext}^n\sK)\right) T^n \text{.}
\end{align*}

\subsection{}\label{III:2-13}

因此, 定理等价于迹公式
\[
  \tr\left(F^*,\eR\Gamma_c(\operatorname{Sym}^n(X,\Gamma_\text{ext}^n \sK)\right) = \sum_{x\in\operatorname{Sym}^n(X)^F} \tr\left(F_x^*,\Gamma_\text{ext}^n\sK\right) \text{.}
\]
在情形 a) 中, 它们由报告 \ref{II:4-10} 得到. 情形 b) 中的类似公式将在 \S\ref{III:4} 证明.

\section{Artin--Schreier 理论}\label{III:3}

\subsection{}\label{III:3-1}

以下总将概形 $S$ 上的凝聚层视为 $S_\text{ét}$ 上的层. 设 $S$ 是特征 $p>0$ 的概形. 若 $G$ 是有限秩 $\dZ/p$-向量空间的局部常值层, 则 $\sG=G\otimes_{\dZ/p}\sO$ 是 $S$ 上的局部自由凝聚层. 我们用 $\Phi$ 同时表示 $\sO$ 的自同态 $f\mapsto f^p$, 及其与 $G$ 的恒等映射的张量积: $\phi:\sG\to\sG$. Artin--Schreier 理论 \cite[IX 3.5]{sga4} 断言下列序列
\[\xymatrix{
  0 \ar[r]
    & \dZ/p \ar[r]
    & \sO \ar[r]^-{\Phi-1}
    & \sO \ar[r]
    & 0
}\]
正合. 与 $G$ 取张量积, 得到正合列
\[\xymatrix{
  0 \ar[r]
    & G \ar[r]
    & \sG \ar[r]^-{\Phi-1}
    & \sG \ar[r]
    & 0
}\]

\subsection{}\label{III:3-2}

假设 $S$ 是 Noether 概形 $\bar S$ 的开子概形. 设 $j:S\hookrightarrow \bar S$ 为包含态射, $I$ 是定义闭补集 $\bar S\setminus S$ 的理想层. 由函子性将 $\Phi$ 延拓到 $j_*\sG$ 所得的映射仍记为 $\Phi$.

\begin{lemma_}\label{III:3-3}
存在 $\sG$ 的凝聚延拓 $\bar\sG\subset j_*\sG$, 使得
\begin{equation*}\tag{3.3.1}\label{III:eq:3-3-1}
  \Phi(\bar\sG)\subset I\cdot \bar\sG\text{.}
\end{equation*}
若 $\bar\sG$ 是这样的延拓, 则序列
\begin{equation*}\tag{3.3.2}\label{III:eq:3-3-2}
\xymatrix{
  0 \ar[r]
    & j_! G \ar[r]
    & \bar\sG \ar[r]^-{\Phi-1}
    & \bar\sG \ar[r]
    & 0
}
\end{equation*}
正合.
\end{lemma_}
\begin{enumerate}[a)]
  \item \emph{存在性}: 设 $\sG'\subset j_*\sG$ 是 $\sG$ 到 $\bar S$ 上的一个凝聚延拓. 对每个 $n$, 有 $\Phi(I^n \sG')\subset I^{p^n}\Phi(\sG')$. 取充分大的 $n$, 使 $I^{n(p-1)-1}\Phi(\sG')\subset \sG'$, 并置 $\bar\sG=I^n\sG'$. 有 $\Phi(\bar\sG)\subset I^{p n}\Phi(\sG')
    =I\cdot I^n\cdot I^{n(p-1)-1}\Phi(\sG')\subset I\cdot I^n\sG'
    = I\bar\sG$.
  \item \emph{$\Phi-1$ 的核}: 在 $S$ 上, 结论由 \ref{III:3-1} 得到; 只须验证 $\ker(\Phi-1)$ 在 $\bar S$ 上平展的 $U$ 上的一个截面 $x$, 在 $\bar S\setminus S$ 的逆像的某个邻域内为零. 事实上, 由于 $x=\Phi(x)$, 由 \eqref{III:eq:3-3-1} 有 $x\in I\bar\sG$; 此外, $x\in I^n\bar\sG\Rightarrow x=\Phi(x)\in\Phi(I^n\bar\sG)\subset I^{n p}\bar\sG$, 所以 $x$ 属于所有 $I^n\bar\sG$, 因而在 $\bar S\setminus S$ 附近为零.
  \item \emph{$\Phi-1$ 的满射性}: 可假设 $\bar S$ 仿射, 并将 $\sG$ 写成 $\sO^n$ 的商; 于是 $\Phi$ 提升为 $p$-线性映射 $\widetilde\Phi:\sO^n\to\sO^n$, 只须验证 $\widetilde\Phi-1$ 为满射:
    \[\xymatrix{
      \sO^n \ar[r]^-{\widetilde \Phi-1} \ar[d]
        & \sO^n \ar[d] \\
      \sG \ar[r]^-{\Phi-1}
        & \sG
    }\]
    层 $\sO^n$ 是群概形 $\dG_a^n$ 的局部截面层, 而 $\widetilde\Phi-1$ 由从 $\dG_a^n$ 到 $\dG_a^n$ 的一个同态诱导. 此同态平展. 要验证它是满射, 只须在每次基变换 $\spec(k)\to\bar S$ ($k$ 代数闭) 后验证; 而在代数闭域 $k$ 上, 同维连通群之间的每个平展同态都是满射. 最后, $\bar S$-概形的平展满射态射在 $\et{\bar S}$ 上对应的层之间诱导满态射.
\end{enumerate}

\subsection{Frobenius 与 Frobenius}\label{III:3-4}

绝对 Frobenius 态射 $F_\text{abs}:S\to S$ 在 $S$ 的底层拓扑空间上为恒等映射, 在结构层上则为 $f\mapsto f^p$. 对平展的 $U/S$, 有典范同构 $F_\text{abs}^* U=U$: 图
\[\xymatrix{
  U \ar[r]^-{F_\text{abs}} \ar[d]
    & S \ar[d] \\
  S \ar[r]^-{F_\text{abs}}
    & S
}\]
为笛卡儿图. 因而, 由 $F_\text{abs}$ 定义的赋环景的自同态 $(\et{\bar S},\sO)\to(\et{\bar S},\sO)$ 可描述如下:
\begin{center}
  $F_\text{abs}^*$ 在 $\et S$ 上为恒等函子; 同态 $F_\text{abs}^*\sO=\sO\to\sO$ 是 $\Phi$.
\end{center}
采用此描述, 对 $S$ 上每个平展层 $\sG$, Frobenius 对应 $F_\text{abs}^*\sG\to\sG$ 就是恒等映射.

给定 $S$ 上拟凝聚层的 $p$-线性态射 $u:\sM\to\sN$, 等价于给定线性态射 $u':F_\text{abs}^*\sM\to\sN$. 特别地, 用 \ref{III:3-5} 的记号, 同态 $\Phi:\sG\to\sG$ 对应于 $\Phi':F_\text{abs}^*\sG\to\sG$.

\begin{lemma_}\label{III:3-5}
沿用 \ref{III:3-4} 的记号, 图
\[\xymatrix{
  F_\textnormal{abs}^* G \ar[r] \ar[d]
    & G \ar[d] \\
  F_\textnormal{abs}^*\sG \ar[r]^-{\Phi'}
    & \sG
}\]
交换.
\end{lemma_}

只须局部验证. 因而可假设 $G$ 常值, 此时引理显然成立.

同一理论也适用于 $F_\text{abs}$ 的幂.

\subsection{}\label{III:3-6}

现在假设 $\bar S$ 是 $\dF_q$ ($q=p^f$) 上的概形, 并依照报告 \S\ref{II:1} 的约定, 用 $S_0,\bar S_0,G_0,\sG_0,\bar\sG_0$ 代替 $S,\dots$. 将引理 \ref{III:3-5} 用于 $F_\text{abs}^F:S_0\to S_0$, 得到交换图
\[\xymatrix{
  F^* G_0 \ar[r] \ar[d]
    & G_0 \ar[d] \\
  F^*\sG_0 \ar[r]^-{(\Phi^f)'} \ar[r]
    & \sG_0
}\qquad\text{, 继而}\qquad
\xymatrix{
  F^*(j_! G_0) \ar[r] \ar[d]
    & j_!G_0 \ar[d] \\
  F^*\bar\sG_0 \ar[r]^-{(\Phi^f)'}
    & \sG_0
}\text{,}
\]
再将标量从 $\dF_q$ 扩张到 $\dF$ 并取上同调, 得
\[\xymatrix{
  F^*(j_! G) \ar[r]^-{F_*} \ar[d]
    & j_!G \ar[d] \\
  F^*\bar\sG \ar[r]
    & \bar\sG
}\qquad\text{及}\qquad
\xymatrix{
  \h^\bullet(\bar S,j_! G) \ar[r]^-{F^*} \ar[d]
    & \h^\bullet(\bar S,j_!G) \ar[d] \\
  \h^\bullet(\bar S,\bar\sG) \ar[r]
    & \h^\bullet(\bar S,\bar\sG)
}\text{,}
\]
其中第二行是将 $\h^\bullet(\bar S_0,\bar G_0)$ 的 $\dF_q$-线性自同态 $\Phi^f$ 从 $\dF_q$ 标量扩张到 $\dF$ 所得的.

Artin--Schreier 理论给出长正合列
\[\xymatrix{
  \cdots \ar[r]
    & \h^i(\bar S,j_! G) \ar[r]
    & \h^i(\bar S,\bar\sG) \ar[r]^-{\Phi-1}
    & \h^i(\bar S,\bar\sG) \ar[r]
    & \cdots
}\]
其中 $\Phi$ 是 $\h^\bullet(\bar S_0,\bar\sG_0)$ 的 $p$-线性自同态 $\Phi$ 到 $\h^\bullet(\bar S,\bar \sG)= \h^\bullet(\bar S_0,\bar\sG_0)\otimes_{\dF_q}\dF$ 的 $p$-线性扩张.

\subsection{}\label{III:3-7}

若 $\bar S_0$ 在 $\dF_q$ 上适当, 则各 $\h^i(\bar S,\bar\sG)$ 有限维, 且 $\Phi-1$ 为满射 (参见 \ref{III:3-6}b). 因此
\[
  \h_c^i(S,G) = \ker\left(\Phi-1 : \h^i(\bar S,\bar G) \to \h^i(\bar S,\bar\sG)\right) \text{,}
\]
而 $F^*$ 由 $\h^i(\bar S_0m\bar\sG_0)$ 的 $\dF_q$-线性自同态 $\Phi^f$ 到 $\h^i(\bar S,\bar\sG) = \h^i(\bar S_0,\bar\sG_0)\otimes_{\dF_q}\dF$ 的 $\dF$-线性扩张诱导.

应用下面的引理 \ref{III:3-8}, 得到等式
\begin{equation*}\tag{3.7.1}\label{III:eq:3-7-1}
  \tr\left(F^*,\h_c^i(S,G)\right) = \tr\left(\Phi^f,\h^i(\bar S_0,\bar\sG_0)\right) \text{.}
\end{equation*}

\begin{lemma_}\label{III:3-8}
设 $\Phi$ 是 $\dF_q$ 上向量空间 $V$ 的一个 $p$-线性自同态: $\phi(\lambda x) = \lambda^p\phi(x)$. 置 $F=\Phi^f$ ($F$ 线性), 并仍用 $\Phi$ 与 $F$ 表示 $\Phi$ 与 $F$ 到 $V\otimes_{\dF_q}\dF$ 的相应 $p$-线性扩张. 则 $F$ 保持 $\ker(\Phi-1)\subset V\otimes_{\dF_q}\dF$ 不变, 且
\[
  \tr\left(F,\ker(\Phi-1)\right) = \tr(F,V) = \tr\left(F,V\otimes_{\dF_q}\dF\right) \text{.}
\]
\end{lemma_}

在 $V\otimes_{\dF_q}\dF$ 上, 有 $F\Phi = \Phi F$: $p$-线性自同态 $F\Phi-\Phi F$ 在 $V$ 上为零, 因而处处为零. 写成 $V=V'+V''$, 其中 $F$ 在 $V'$ 上可逆, 在 $V''$ 上幂零. $p$-线性自同态理论保证
\[
  \ker(\Phi-1)\otimes_{\dF_p}\dF \iso V'\otimes_{\dF_q}\dF \text{.}
\]
由于 $\tr(F,V'') = 0$, 因而有
\begin{align*}
  \tr\left(F,\ker(\Phi-1)\right)
    &= \tr\left(F,\ker(\Phi-1)\otimes_{\dF_q}\dF\right) = \tr\left(F,V'\otimes_{\dF_q}\dF\right) \\
    &= \tr\left(F,V\otimes_{\dF_q}\dF\right) = \tr(F,V) \text{.}
\end{align*}

\section{模 \texorpdfstring{$p$}{p} 迹公式}\label{III:4}

\begin{theorem_}\label{III:4-1}
设 $X_0$ 是 $\dF_q$ 上的分离有限型概形, $A$ 是特征为 $p$ 的交换既约 Noether 环, $\sK_0\in\ob\eD_\textnormal{ctf}(X_0,A)$. 有
\begin{equation*}\tag{4.1.1}\label{III:eq:4-1-1}
  \sum_{x\in X^F} \tr\left(F_x^*,\sK_x\right) = \tr\left(F^*,\eR\Gamma_c(X,\sK)\right) \text{.}
\end{equation*}
\end{theorem_}

\subsection{归约到 \texorpdfstring{$A$}{A} 为有限域的情形}\label{III:4-2}

% NOTE: 英文源指出本段及其他段中的 (3.1.1) 似乎指 (4.1.1).
由标准的取极限论证, 可归结到 $A$ 在 $\dZ/p$ 上为有限型的情形. 对 $A$ 的每个极大理想 $\fm$, \eqref{III:eq:4-1-1} 两边在 $A/\fm$ 中的像都由同一公式给出, 其中将 $\sK_0$ 换为 $\sK_0\lotimes_A A/\fm\in\ob\eD_\text{ctf}(X_0,A/\fm)$ (右边的情形参见报告 \ref{II:4-12}). 由于各 $A/\fm$ 都有限, 而 $A$ 的各极大理想之交为 $0$, 即得到所述归约.

\subsection{归约到 \texorpdfstring{$A=\dZ/p$}{A=Z/p} 的情形}\label{III:4-3}

假设 $A$ 为有限域, 用 $T_A'(\sK_0)$ 和 $T_A''(\sK_0)$ 表示 \eqref{III:eq:4-1-1} 的两边. 若通过限制标量, 将 $\sK_0$ 看作 $\dZ/p$-模层的复形, 则有
\[
  T_{\dZ/p}'(\sK_0) = \tr_{A/(\dZ/p)} T_A'(\sK_0)
\]
对 $T''$ 也一样. 因而, 若 \eqref{III:eq:4-1-1} 对 $A=\dZ/p$ 成立, 则有
\begin{equation*}\tag{4.3.1}\label{III:eq:4-3-1}
  \tr_{A/(\dZ/p)} T_A'(\sK_0) = \tr_{A/(\dZ/p)} T_A''(\sK_0) \text{.}
\end{equation*}
我断言, 对每个 $\lambda\in A$, 甚至有
\begin{equation*}\tag{4.3.2}\label{III:eq:4-3-2}
  \tr_{A/(\dZ/p)}\left(\lambda\cdot T_A'(\sK_0)\right) = \tr_{A/(\dZ/p)}\left(\lambda\cdot T_A''(\sK_0)\right) \text{,}
\end{equation*}
从而 $T_A'(\sK_0) = T_A''(\sK_0)$. 当 $\lambda=0$ 时这是显然的. 当 $\lambda$ 可逆时, 以 $A(\lambda)$ 表示 $\spec(\dF_q)$ 上满足 $F^*=\lambda$ 的秩一自由 $A$-模层在 $X$ 上的逆像. 公式 \eqref{III:eq:4-3-2} 恰好是将 \eqref{III:eq:4-3-1} 用于 $A(\lambda)\otimes_A \sK_0$ 所得的公式.

% 原编号为定理 4.1bis (即再次陈述定理 4.1).
\begin{theorem_}\label{III:4-1bis}
设 $\bar S_0$ 是光滑射影曲线, $j:S_0\hookrightarrow \bar S_0$ 是 $\bar S_0$ 的稠密开子概形, $G_0$ 是 $S_0$ 上有限秩 $\dZ/p$-向量空间的局部常值层. 有
\begin{equation*}\tag{4.4.1}\label{III:eq:4-4-1}
  \sum_{x\in S^F} \tr\left(F_x^*,G_x\right) = \sum_i (-1)^i \tr\left(F^*,\h_c^i(S,G)\right) \text{.}
\end{equation*}
\end{theorem_}
\begin{proof}
将 \ref{III:3-3} 应用于 $(S_0,\bar S_0,G_0)$. 由 \ref{III:3-3} 保证存在的 $\sG_0$ 的凝聚延拓 $\bar\sG_0$ 是局部自由的, 因为它无挠 (有 $\bar\sG_0\subset j_*\sG_0$), 而 $\bar S_0$ 是一维正则概形. $\bar\sG_0$ 的 $q$-线性自同态 $\Phi^f$ 可解释为态射
\[
  u : F^*\bar\sG_0 \to \bar\sG_0 \text{.}
\]
它通过 $I\bar G$ 分解, 其中 $I$ 定义闭子集 $\bar S\setminus S$. 令 $v:j_!G_0 \to \bar\sG_0$ 为自然映射 $G_0 \to \sG_0 = G_0\otimes_{\dZ/p}\sO$ 的唯一延拓. 图为
\begin{equation*}\tag{4.4.2}\label{III:eq:4-4-2}
\xymatrix{
  F^* j_! G_0 \ar[r]^-{F^*} \ar[d]^-{F^* v}
    & j_! G_0 \ar[d]^-v \\
  F^*\bar \sG_0 \ar[r]^-u
    & \bar\sG_0
}
\end{equation*}
最后, 用 $u F^*$ 表示复合映射
\[\xymatrix{
  \h^\bullet(\bar S,\bar\sG) \ar[r]^-{F^*}
    & \h^\bullet(\bar S,F^*\bar\sG) \ar[r]^-u
    & \h^\bullet(\bar S,\bar\sG) \text{,}
}\]
则由 \eqref{III:eq:3-7-1}, 有
\begin{equation*}\tag{4.4.3}\label{III:eq:4-4-3}
  \tr\left(F^*,\h_c^i(S,G)\right) = \tr\left(u F^*,\h^i(\bar S,\bar\sG)\right) \text{.}
\end{equation*}

% NOTE: 英文源提示须查明 Woods Hole 讨论班的参考文献.
对 $\bar S$ 的每个闭点 $i_x:x\hookrightarrow \bar S$, 置 $\bar\sG_x = i_x^*\bar\sG$. 凝聚上同调中的 Lefschetz 迹公式 (“Woods Hole fixed point formula”), 应用于 $F$ 和 $u$, 说明
\begin{equation*}\tag{4.4.4}\label{III:eq:4-4-4}
  \sum_i (-1)^i \tr\left(uF^*, \h^i(\bar S,\bar\sG)\right) = \sum_{x\in \bar S^F} \frac{\tr(u_x,\sG_x)}{\det(1-d F_x)} \text{.}
\end{equation*}
在右边,
\begin{enumerate}[\indent a)]
  \item 由于 $d F=0$, 各分母等于 1;
  \item 对 $x\in \bar S^F\setminus S^F$, $\bar\sG_x$ 的自同态 $u_x$ 为零, 因而 $\tr(u_x,\bar\sG_x)$ 也为零;
  \item 对 $x\in S^F$, \eqref{III:eq:4-4-4} 给出的贡献是
    \[
      \tr(u_x,\bar\sG_x) = \tr(F_x^*,G_x) \text{.}
    \]
    因此, 等式 \eqref{III:eq:4-4-1} 由 \eqref{III:eq:4-4-3} 及 \eqref{III:eq:4-4-4} 得到.
\end{enumerate}
\end{proof}

\subsection{一个反例}\label{III:4-5}

下面的例子说明, \ref{III:4-1} 中 $A$ 既约的假设不可省略.

设 $Y_0$ 是 $\dF_2$ 上 Frobenius 特征值为 $\pm \sqrt{-2}$ 的椭圆曲线. 它有 $3$ 个有理点, 且是超奇异的. 因而对合 $\sigma:x\mapsto -x$ 只有一个不动点, 即点 $0$. $Y_0$ 关于 $\sigma$ 的商是一条射影直线; 因此 $Y_0\setminus 0$ 关于 $\sigma$ 的商是一条仿射直线 $\dF_0^1$, 且 $Y_0\setminus 0$ 是 $\dF_0^1$ 的无分歧二重覆盖:
\[
  \pi:Y_0\setminus 0 \to \dA_0^1 \qquad \text{, $\pi=\pi\sigma$.}
\]

层 $\pi_*\dZ/2$ 是 $\dZ/2$ 上二元素群 $\{1,\sigma\}$ 的群代数 $A$ 上的秩一自由模. 此代数同构于 $\dZ/2$ 上的数的代数.

$A$-模层 $\pi_*\dZ/2$ 就是所要的例子:
\begin{enumerate}[\indent a)]
  \item 有 $\h_c^\bullet(\dA^1,\pi_*\dZ/2) = \h_c^\bullet(Y\setminus 0,\dZ/2)$. 由于 $Y_0$ 超奇异, 当 $i>0$ 时有 $\h^i(Y,\dZ/2)=0$, 而长正合列
    \[
      \cdots \to \h_c^i(Y\setminus 0,\dZ/2) \to \h^i(Y,\dZ/2) \to \h^i(\{0\},\dZ/2) \to \cdots
    \]
    说明 $\h_c^\bullet(Y\setminus 0,\dZ/2)=0$. 因此 $\h_c^\bullet(\dA^1,\pi_*\dZ/2)=0$, 且
    \[
      \tr_A\left(F^*,\eR\Gamma_c(\dA^1,\pi_*\dZ/2)\right) = 0 \text{.}
    \]
  \item 由于 $Y_0\setminus 0$ 有两个有理点, 在 $\dA_1^0$ 的一个有理点处, Frobenius 置换为恒等映射, 而在另一个有理点处为 $\sigma$. 因此
    \[
      \sum_{x\in {\dA^1}^F} \tr\left(F^*,\pi_*\dZ/2\right) = 1+\sigma \text{,}
    \]
    且 $1+\sigma\ne 0$.
\end{enumerate}

覆盖 $\pi$ 的方程为 $y^2-y=x^3$.
